\section{Beyond Gaussian sources: a radial confinement obstruction}
\label{app:radial-extension}
The source two-jet is the pair of first and second time derivatives of
the distribution path at its start. The same calculation can be made for
some non-Gaussian sources. To complete the obstruction proof, we must
bound the squared correction perpendicular to a selected direction.
We prove that bound for the radial source class below. Odd dimension alone
would not supply the needed moment estimate.\AIassistance

Let
\begin{equation}
 q_W(x)=Z_W^{-1}\exp\{-|x|^2/2-W(|x|^2)\},
 \quad W\in C^2([0,\infty)),\quad W(0)=0,\quad W'\ge0.
 \label{eq:radial-source}
\end{equation}
The normalizing constant makes the density integrate to one. The
additional radial term reduces the relative weight of large radii.
This includes $W(s)=\epsilon s^2/4$, $\epsilon>0$, which is non-Gaussian.
The density is positive and smooth, and its moments of every order are finite. Let
$X\sim q_W$ and $Y$ be independent, with $\E Y=0$, $Y\in\mathcal P_2$,
and covariance $S\succ0$. Prescribe $p_t=\Law((1-t)X+tY)$ near zero.

\begin{theorem}[Non-Gaussian source-local obstruction]
\label{thm:radial-proof}
If $d$ is odd, this path has no $C^1$ generator near zero that is globally
Lipschitz in space, whose source acceleration satisfies $a_0\in L^1(q_W)$,
and whose source acceleration preserves the source distribution:
$\operatorname{div}(q_Wa_0)=0$. Every field in $\overline{\mathcal B}_K$
satisfies the integrability condition. Differentiation of compact flow observables
twice at zero must be valid; the common derivative bound in the paper is a sufficient,
but not necessary, condition. No Gaussian, symmetry, or smooth-density
assumption is made on the target.
\end{theorem}
\begin{proof}
The proof follows the Gaussian argument until the moment estimate.
The first path derivative says that the correction leaves the source
distribution unchanged. The second derivative relates the squared
correction in a direction to acceleration in that direction.
We then bound moments in the perpendicular directions for this source.
With those bounds, the large-sphere argument again produces a test that
detects nonzero acceleration.

Put $r=u_0+x$ and $L=\operatorname{Lip}(r)$. Differentiating the path once
gives $\operatorname{div}(q_Wr)=0$. Since $r$ is $C^1$, this identity is
pointwise and gives
\[
 \operatorname{div}r=(1+2W'(|x|^2))x\cdot r,
 \qquad |x\cdot r(x)|\le dL.
\]
For a unit $e$, the marginal flux
$q_{W,e}(y)\E[e\cdot r(X)\mid e\cdot X=y]$ has zero distributional
derivative. It is integrable and hence identically zero. The source
second-jet argument consequently gives exactly
\begin{equation}
 \E[\psi'(e\cdot X)e\cdot a_0(X)]
 =\int\psi''(y)\{e^\top Se-Q_e(y)\}q_{W,e}(y)\,dy,
 \quad Q_e(y)=\E[(e\cdot r(X))^2\mid e\cdot X=y].
 \label{eq:radial-twojet}
\end{equation}
Indeed $(e\cdot(Y-X))^2$ has conditional mean $y^2+e^\top Se$, while
$(e\cdot u_0(X))^2$ has conditional mean $y^2+Q_e(y)$; the cross term
vanishes by the flux identity. Cutoffs, finite second moments, and
$a_0\in L^1(q_W)$ justify the extension from compact tests to the scalar
bounded-derivative tests used below: the cutoff error is at most a constant
times $\int_{|x|\ge R}|a_0|q_W\dd x$.

Conditional on $e\cdot X=y$, the perpendicular coordinate has density
proportional to
$\exp(-|z|^2/2)\exp\{-W(y^2+|z|^2)\}$.
Relative to a standard Gaussian perpendicular coordinate, this density
gives smaller relative weight to larger radii. The weight is nonincreasing
because $W'\ge0$. The following elementary comparison therefore bounds
its increasing radial moments by the Gaussian moments. For an increasing function $f$ and a
decreasing nonnegative $w$, $\E[f(R)w(R)]\le\E f(R)\E w(R)$: subtract
using two independent copies of $R$ and note that
$(f(R)-f(R'))(w(R)-w(R'))\le0$.
Thus the conditional moments satisfy
\[
 \E|Z_\perp|^2\le d-1,\qquad
 \E|Z_\perp|^4\le(d-1)(d+1).
\]
The same argument shows finiteness of all needed source moments.

Let $A_d=d+\sqrt{d^2+d-2}$ and choose
$R\ge\max\{1,2A_dL/\sqrt{\lambda_{\min}(S)}\}$.
The tangential projection of $r$ has a zero at $Re$ for some $e$ by the
hairy-ball theorem on $S^{d-1}$; for $d=1$ its tangent space is zero.
For $y\in[R,R+1]$, the pointwise divergence identity and Lipschitzness give
\[
 |e\cdot r(ye+z)|\le\frac L R
       \{d+|z|\sqrt{(y-R)^2+|z|^2}\}.
\]
The conditional moment bounds therefore imply
$\sqrt{Q_e(y)}\le A_dL/R\le\sqrt\lambda/2$.
A convex scalar test with $\psi''=\mathbf1_{[R,R+1]}$ and
$0\le\psi'\le1$ has positive pairing at least
$(3\lambda/4)\Pr\{e\cdot X\in[R,R+1]\}>0$ in
\eqref{eq:radial-twojet}. This contradicts source acceleration preservation.
An explicit positive pairing bound follows by bounding the probability below
by
\begin{equation}
 \frac{\operatorname{vol}(B_{d-1}(1))}{(2\pi)^{d/2}}
 \exp\left\{-\frac{(R+1)^2+1}{2}
              -W((R+1)^2+1)\right\},
 \label{eq:radial-strip-mass}
\end{equation}
using $Z_W\le(2\pi)^{d/2}$ and integrating over the perpendicular unit
ball. For $d=1$ use $\operatorname{vol}(B_0)=1$; the displayed bound is
still conservative.
\end{proof}

The two-grid characterization also allows $\gamma_d$ to be replaced by
a fixed positive continuous source density with a finite fourth moment.
The changes to the proof are specific. In
Lemmas~\ref{lem:compact}--\ref{lem:weak}, use the source's second and
fourth moments in place of Gaussian moment bounds. Second-moment tails
justify compactness and necessity. The fourth moment controls the
acceleration-flow coupling used for sufficiency, with physical time fixed.
A $C^1$ bi-Lipschitz flow preserves the positive continuous density needed
in that proof. The preceding
obstruction therefore gives a uniform target-only $c/N$ barrier for
\eqref{eq:radial-source} under the same common derivative bound and balanced
refinements. The quantitative Hermite result does \emph{not} extend to
this non-Gaussian source without a corresponding spectral argument.

An invertible affine change of source coordinates preserves the source
obstruction and Euler update algebra. If $X=A\widetilde X+b$, transform
velocities and target covariances by $A^{-1}$; the derivative bound changes by the chain rule, including the affine
shift and the operator norms of $A,A^{-1}$, and
$W_2(A_\#\mu,A_\#\nu)\ge s_{\min}(A)W_2(\mu,\nu)$.
This covers non-isotropic elliptical versions of the source class. It does
not cover arbitrary multimodal initial densities, arbitrary heavy tails,
or discontinuous velocities. It also does not cover velocities assumed only
to have weak derivatives in a Sobolev space.

\section{The exact asymptotic optimum over all scalar nonlinear candidates}
\label{app:sharp-minimax}
The characterization gives the error's order of decay, but it generally
does not give the smallest coefficient multiplying that order.
For the scalar Gaussian path below, we find that coefficient over
\emph{all} admissible nonlinear fields. Each field is judged by the sum of
its worst coarse and refined errors over the tested times. We then choose
the best field. This order of worst-time and best-field choices defines the
minimax problem. No general sharp constant is established above one dimension.
Let $d=T=1$ and $p_t=\mathcal N(0,g(t)^2)$ with
$g(t)=\sqrt{(1-t)^2+t^2}$. Use uniform $N$-step Euler and its fixed
$m$-fold refinement, $m\ge2$. Write $q_N^{(m)}=\inf_{u\in\mathcal B_K}
(E_c+E_r)$.\AIassistance

\begin{theorem}[Sharp all-candidate scalar constant]
\label{thm:sharp-minimax-proof}
For any fixed derivative bound strictly exceeding
$\max\{1,\|c\|_\infty,\|c'\|_\infty,\|c''\|_\infty\}$,
where $c=g'/g$ on $[0,1]$, the secant fields belong to $\mathcal B_K$ for
all sufficiently large $N$, and
\begin{equation}
 \lim_{N\to\infty}Nq_N^{(m)}
 =\frac12(1-1/m)(1+\pi/2).
 \label{eq:sharp-minimax}
\end{equation}
In particular the bisection constant is $(2+\pi)/8$.
\end{theorem}
\begin{proof}
We first show that an asymptotically accurate sequence must approach the
unique exact scalar generator. A uniform particle-error expansion then gives
a lower bound on the separation between its coarse and fine numerical laws.
Finally, the secant field attains that lower bound.

Any globally Lipschitz exact generator of this scalar Gaussian path is
$v(t,x)=c(t)x$. Indeed the difference of two generators satisfies
$\partial_x(p_tr_t)=0$; Gaussian integrability and at most linear growth
force its constant flux to be zero. This is the standard one-dimensional uniqueness argument for a generator
with vanishing flux.

We first record an elementary uniform particle expansion on uniform grids.
For any $u\in\overline{\mathcal B}_K$, with $X_t=F^u_{0,t}(G)$, let
\[
 B_t^u=J_t(G)\int_0^tJ_s(G)^{-1}a_s(X_s)\,ds.
\]
The common-start recursive Euler error satisfies, in $L^2$ uniformly over
nodes and the derivative class,
\begin{equation}
 \widehat X_k-X_{kh}=-\tfrac h2 B_{kh}^u+O_{L^2}(h^2).
 \label{eq:scalar-variational-expansion}
\end{equation}
To derive the expansion, write $e_k=\widehat X_k-X_{kh}$. Taylor's
formula gives
$e_{k+1}=(1+hDu(t_k,X_{t_k}))e_k-h^2a_{t_k}(X_{t_k})/2+r_k$,
where $|r_k|\le C h^3(1+|G|)^2$ after using the uniform strong bound
$|e_k|\le Ch(1+|G|)$ and $\|D^2u\|\le K$.
Variation of constants for $B'=Du(t,X_t)B+a(t,X_t)$ has a one-step
quadrature remainder $O_{L^2}(h^2)$; multiplying it by $h/2$ and comparing
the recurrences yields \eqref{eq:scalar-variational-expansion} by discrete
Gronwall. Gaussian fourth moments bound all remainder terms. This uses
only the derivatives already controlled by the class.

Choose approximate minimizers with $E_c+E_r=O(1/N)$, possible by the
secant upper construction below. Every convergent subsequence has, by the
node-to-path bound and compactness, an exact limiting generator and hence
limit $v$. Local $C^1$ convergence plus Gaussian tails implies
$B^{u_N}\to B^v$ in uniform-in-time $L^2$. Subtract
\eqref{eq:scalar-variational-expansion} for steps $h$ and $h/m$ to obtain
\[
 h^{-1}(\widehat X^c_k-\widehat X^r_{mk})
 \longrightarrow-\tfrac12(1-1/m)B^v_{t_k}.
\]
For $hK<1$, each scalar Euler update is increasing. Particles therefore
keep the order of their initial positions. Pairing the same initial
quantiles also pairs the same output quantiles. Such a pairing is
\emph{optimal} in one dimension. Thus the $W_2$ distance between the two
numerical laws equals the displayed common-start $L^2$ norm. The target triangle inequality proves the lower bound with
constant $\tfrac12(1-1/m)\sup_t\|B_t^v\|_2$.

For $v=cx$, direct integration gives
$B_t^v=g(t)\int_0^t g''(s)/g(s)\,ds\,G$.
Let $I(t)=\int_0^t1/g(s)^4\,ds$. The product $g(t)I(t)$ is increasing.
On $t\le1/2$, $I(t)\le t/g(t)^4$, so its derivative is at least
$(1-3t+4t^2)/g(t)^5>0$; on $t\ge1/2$ both terms in the derivative are
nonnegative. Its maximum is thus $I(1)=1+\pi/2$, obtained by substituting
$z=2s-1$.

Finally $u_h(t,x)=[g(t+h)-g(t)]x/[hg(t)]$ converges to $v$ in the scalar
coefficient $C^2$ norm and satisfies the derivative bound strictly for sufficiently small
$h$. It gives exact coarse target marginals, while its refined target
error has coefficient $\tfrac12(1-1/m)\sup_t g(t)I(t)$. This supplies the
matching upper bound and proves the limit. The optimal constant concerns
uniform grids, this fixed scalar path, and the sum-of-maxima objective;
it is not a sharp constant for general dimension or optimized nonuniform
grids.
\end{proof}

\section{One-dimensional geometry and the limit of endpoint tests}
\label{app:observable-geometry}
A quantile function gives the position below which a specified proportion
of a distribution lies. For example, the middle quantile has half the
probability below it. In one dimension, following each fixed quantile
through time exactly transports the distribution. We can therefore read
the acceleration condition directly from these quantile trajectories.\AIassistance

Let $p_t$ be a scalar path with a positive continuous density and a finite
second moment, and let $Q_t$ be its quantile function. Among the fields in
$\overline{\mathcal B}_K$, the flux argument of Supplementary
Section~\ref{app:sharp-minimax} shows that the unique exact generator is
$u(t,Q_t(z))=\partial_tQ_t(z)$ and
$a(t,Q_t(z))=\partial_{tt}Q_t(z)$.
Thus an order-elevating generator exists exactly when
$\partial_{tt}Q_t=0$ in $L^2(0,1)$ and the resulting generator belongs to
the chosen class closure. Necessity follows because the scalar weighted flux must vanish.
For sufficiency, each quantile moves along a straight trajectory, so Euler
reproduces that motion exactly. This is the standard one-dimensional
displacement-interpolation criterion. It constructs a field for an
externally specified scalar path. In several dimensions, a straight,
invertible interpolation of prescribed transport maps also gives a
sufficient Euler-exact construction. This construction is known in
displacement geometry and is not necessary in general.

Accuracy only at the final time cannot replace accuracy along the whole
path. Errors with different signs can cancel before the final time. The smooth
scalar field $u(t,x)=-5t^2x$ generates $g(t)=\exp(-5t^3/3)$, and
\[
 \int_0^1\{c'(s)+c(s)^2\}\,ds
 =\int_0^1(-10s+25s^4)\,ds=0.
\]
The terminal Euler error is $O(h^2)$ simultaneously for every fixed
integer refinement factor, while the acceleration is nonzero and the
maximum-node error is first-order. Conversely, under a \emph{specified}
transported test whose pairing has one sign and satisfies $|I(t)|\ge a>0$
on the full interval, \eqref{eq:balanced-contrast} with $i=0,j=N$ and
$Z_0=0$ gives a terminal $c/N$ bound whenever the remainder is dominated.
The extra sign condition is an observability assumption; it is not implied
by a nonzero local residual.
